\documentclass[a4paper,12pt]{article}

% 页面设置
\usepackage[top=2.54cm, bottom=2.54cm, left=1.91cm, right=1.91cm]{geometry}

% 中文支持
\usepackage[UTF8]{ctex}
\usepackage{fontspec}

% 数学公式支持
\usepackage{amsmath}
\usepackage{amssymb}
\usepackage{amsfonts}
\usepackage{amsthm}

\usepackage{enumitem}

% 定理环境定义
\theoremstyle{definition}
\newtheorem*{pf}{证}
\newtheorem*{sol}{解}

% 图形支持
\usepackage{graphicx}
\usepackage{tikz}

% 超链接支持
\usepackage{hyperref}
\hypersetup{
    colorlinks=true,
    linkcolor=blue,
    filecolor=magenta,
    urlcolor=cyan,
}

% 行距设置
\linespread{1.25}

% 标题信息
\title{Mathematical Analysis II\\Week 9 Homework (April 30)}
\author{袁子强\hspace{1cm} 2025533009}
% \date{}

\begin{document}

\maketitle

\section*{Section 10.3}
15. 求由曲面 $\displaystyle\frac{x^2}{a^2} + \frac{y^2}{b^2} = \frac{z^2}{c^2}$ 和 $z = c$ 所围成的均匀物体的重心坐标.
\begin{sol}
    由于区域关于 $z$ 轴旋转对称且密度均匀,因此仅需计算 $z$ 坐标 $\dfrac{\iiint_V z\,\mathrm{d}V}{\iiint_V\,\mathrm{d}V}$.

    采用广义柱面坐标变换: 记
    $$\begin{cases}
        x=ar\cos\theta \\
        y=br\sin\theta \\
        z=z
    \end{cases},$$
    则积分区域为 $\begin{cases}
        0\leq r \leq \dfrac{z}{c} \\
        0\leq \theta \leq 2\pi
    \end{cases}$, 雅可比行列式为 $\dfrac{\partial(x,y,z)}{\partial(r,\theta,z)}=abr$.

    因此

    \begin{align*}
        \iiint_V z\,\mathrm{d}V
         & =ab\int_0^c z\,\mathrm{d}z\int_0^{2\pi}\,\mathrm{d}\theta\int_0^{\frac{z}{c}} r\,\mathrm{d}r \\
         & =\dfrac{ab\pi}{c^2}\int_0^c z^3\,\mathrm{d}z                                                  \\
         & =\dfrac{abc^2\pi}{4}
    \end{align*}

    \begin{align*}
        \iiint_V \,\mathrm{d}V
         & =ab\int_0^c \,\mathrm{d}z\int_0^{2\pi}\,\mathrm{d}\theta\int_0^{\frac{z}{c}} r\,\mathrm{d}r \\
         & =\dfrac{ab\pi}{c^2}\int_0^c z^2\,\mathrm{d}z                                                 \\
         & =\dfrac{abc\pi}{3}
    \end{align*}

    从而 $z$ 坐标为 $\dfrac{\frac{abc^2\pi}{4}}{\frac{abc\pi}{3}}=\dfrac{3c}{4}$.

    因此重心坐标为 $\left(0,0,\dfrac{3c}{4}\right)$.
\end{sol}

16. 设球体 $x^2 + y^2 + z^2 \leq 2az$ 内各点密度与各点到原点的距离成反比,求其重心坐标.
\begin{sol}
    原方程可化为 $x^2 + y^2 + (z-a)^2 \leq a^2$,表示球心为 $(0,0,a)$、半径为 $a$ 的球体.由于区域关于 $z$ 轴旋转对称,因此仅需计算 $z$ 轴坐标.设密度与到原点距离的比例系数为 $k$,即点 $(x,y,z)$ 处的密度为 $\dfrac{k}{\sqrt{x^2+y^2+z^2}}$.

    采用球面坐标变换: 记
    $$\begin{cases}
        x=r\sin\theta\cos\varphi \\
        y=r\sin\theta\sin\varphi \\
        z=r\cos\theta
    \end{cases},$$
    则积分区域为 $\begin{cases}
        0\leq r \leq 2a\cos\theta       \\
        0\leq \theta \leq \dfrac{\pi}{2} \\
        0\leq \varphi \leq 2\pi
    \end{cases}$, 雅可比行列式为 $\dfrac{\partial(x,y,z)}{\partial(r,\theta,\varphi)}=r^2\sin\theta$.

    因此

    \begin{align*}
        \iiint_V z\rho\left(\overrightarrow{r}\right)\,\mathrm{d}V
         & =k\int_0^{\frac{\pi}{2}}\sin\theta\cos\theta\,\mathrm{d}\theta\int_0^{2\pi}\,\mathrm{d}\varphi\int_0^{2a\cos\theta}r^2\,\mathrm{d}r \\
         & =\dfrac{16a^3k\pi}{3}\int_0^{\frac{\pi}{2}}\sin\theta\cos^4\theta\,\mathrm{d}\theta                                                  \\
         & =\dfrac{16a^3k\pi}{15}
    \end{align*}

    \begin{align*}
        \iiint_V \rho\left(\overrightarrow{r}\right)\,\mathrm{d}V
         & =k\int_0^{\frac{\pi}{2}}\sin\theta\,\mathrm{d}\theta\int_0^{2\pi}\,\mathrm{d}\varphi\int_0^{2a\cos\theta}r\,\mathrm{d}r \\
         & =4a^2k\pi\int_0^{\frac{\pi}{2}}\sin\theta\cos^2\theta\,\mathrm{d}\theta                                                 \\
         & =\dfrac{4a^2k\pi}{3}
    \end{align*}

    从而 $z$ 轴坐标为 $\dfrac{\frac{16a^3k\pi}{15}}{\frac{4a^2k\pi}{3}}=\dfrac{4a}{5}$.

    因此重心坐标为 $\left(0,0,\dfrac{4a}{5}\right)$.
\end{sol}

18. 求以下各物体的转动惯量 (设 $\rho = $ 常数)

(1) 质量为 $m$, 半径为 $R$ 的薄圆盘, 对于:

(a) 通过圆心并垂直圆盘的轴;
\begin{sol}
    取 $z$ 轴为转轴.

    需计算 $\iint_D d^2\rho\,\mathrm{d}A$.

    采用极坐标变换: 记
    $$\begin{cases}
        x=r\cos\theta \\
        y=r\sin\theta
    \end{cases},$$
    则积分区域为 $\begin{cases}
        0\leq r \leq R \\
        0\leq \theta \leq 2\pi
    \end{cases}$, 雅可比行列式为 $\dfrac{\partial(x,y)}{\partial(r,\theta)}=r$, 且 $d^2=r^2$.

    面密度 $\rho=\dfrac{m}{\pi R^2}$.

    因此

    \begin{align*}
        \iint_D d^2\rho\,\mathrm{d}A
         & =\dfrac{m}{\pi R^2}\int_0^{2\pi}\,\mathrm{d}\theta\int_0^R r^3\,\mathrm{d}r \\
         & =\dfrac{mR^2}{2}
    \end{align*}
\end{sol}

(b) 直径.
\begin{sol}
    取 $x$ 轴为转轴.

    此时 $d^2=y^2=r^2\sin^2\theta$.

    因此

    \begin{align*}
        \iint_D d^2\rho\,\mathrm{d}A
         & =\dfrac{m}{\pi R^2}\int_0^{2\pi}\sin^2\theta\,\mathrm{d}\theta\int_0^R r^3\,\mathrm{d}r \\
         & =\dfrac{mR^2}{4}
    \end{align*}
\end{sol}

(2) 质量为 $m$, 半径为 $R$ 的球体, 对于通过球心的轴线.
\begin{sol}
    取 $z$ 轴为转轴,需计算 $\iiint_V d^2\rho\,\mathrm{d}V$.

    采用球面坐标变换: 记
    $$\begin{cases}
        x=r\sin\theta\cos\varphi \\
        y=r\sin\theta\sin\varphi \\
        z=r\cos\theta
    \end{cases},$$
    则积分区域为 $\begin{cases}
        0\leq r \leq R        \\
        0\leq \theta \leq \pi \\
        0\leq \varphi \leq 2\pi
    \end{cases}$, 雅可比行列式为 $\dfrac{\partial(x,y,z)}{\partial(r,\theta,\varphi)}=r^2\sin\theta$, 且 $d^2=r^2\sin^2\theta$.

    体密度 $\rho=\dfrac{m}{\frac{4}{3}\pi R^3}=\dfrac{3m}{4\pi R^3}$.

    因此

    \begin{align*}
        \iiint_V d^2\rho\,\mathrm{d}V
         & =\dfrac{3m}{4\pi R^3}\int_0^{2\pi}\,\mathrm{d}\varphi\int_0^{\pi}\sin^3\theta\,\mathrm{d}\theta\int_0^R r^4\,\mathrm{d}r \\
         & =\dfrac{3m}{4\pi R^3}\cdot 2\pi\cdot\dfrac{4}{3}\cdot\dfrac{R^5}{5}                                                        \\
         & =\dfrac{2}{5}mR^2
    \end{align*}
\end{sol}

19. 求密度为 $\rho$ 的均匀球锥体对于在其顶点为一单位质量的质点的吸引力, 设球的半径为 $R$, 而轴截面的扇形的角等于 $2\alpha$.
\begin{sol}
    取 $z$ 轴为对称轴,顶点为原点.由于垂直方向上的引力分量相互抵消,因此仅需计算沿 $z$ 轴方向的引力分量.

    采用球面坐标变换: 记
    $$\begin{cases}
        x=r\sin\theta\cos\varphi \\
        y=r\sin\theta\sin\varphi \\
        z=r\cos\theta
    \end{cases},$$
    则积分区域为 $\begin{cases}
        0\leq r \leq R           \\
        0\leq \theta \leq \alpha \\
        0\leq \varphi \leq 2\pi
    \end{cases}$, 雅可比行列式为 $\dfrac{\partial(x,y,z)}{\partial(r,\theta,\varphi)}=r^2\sin\theta$.位置 $(r,\theta,\varphi)$ 处质量元对质点的引力在 $z$ 轴方向的分量为 $\dfrac{G\rho}{r^2}\cos\theta$.

    因此

    \begin{align*}
        \iiint_V\dfrac{G\rho}{r^2}\cos\theta\,\mathrm{d}V
         & =G\rho\int_0^{2\pi}\,\mathrm{d}\varphi\int_0^{\alpha}\sin\theta\cos\theta\,\mathrm{d}\theta\int_0^R\,\mathrm{d}r \\
         & =\pi RG\rho\sin^2\alpha
    \end{align*}
\end{sol}

\section*{Section 10.4}

1. 计算下列 $n$ 重积分.

(1) $\displaystyle\int \cdots \int_{[0,1]^n} (x_1^2 + \cdots + x_n^2) \, \mathrm{d}x_1 \cdots \mathrm{d}x_n$;
\begin{sol}
    首先计算 $\int \cdots \int_{[0,1]^n} x_i^2 \, \mathrm{d}x_1 \cdots \mathrm{d}x_n$ ($i=1,2,\dots,n$):

    \begin{align*}
        \int \cdots \int_{[0,1]^n} x_i^2 \, \mathrm{d}x_1 \cdots \mathrm{d}x_n
         & =\int_0^1\,\mathrm{d}x_1\cdots\int_0^1 x_i^2\,\mathrm{d}x_i\cdots\int_0^1\,\mathrm{d}x_n \\
         & =1\cdot\dfrac{1}{3}\cdot 1                                                                \\
         & =\dfrac{1}{3}
    \end{align*}

    可见,各部分单独积分结果与 $i$ 无关.

    因此 $\int \cdots \int_{[0,1]^n} (x_1^2 + \cdots + x_n^2) \, \mathrm{d}x_1 \cdots \mathrm{d}x_n=\dfrac{n}{3}$.
\end{sol}

(2) $\displaystyle\int \cdots \int_{[0,1]^n} (x_1 + \cdots + x_n)^2 \, \mathrm{d}x_1 \cdots \mathrm{d}x_n$;
\begin{sol}
    展开得 $(x_1 + \cdots + x_n)^2=\sum x_i^2+2\sum_{i\neq j}x_ix_j$.

    对于 $n$ 个 $x_i^2$ 项,由(1)可得 $\int \cdots \int_{[0,1]^n} (x_1^2 + \cdots + x_n^2) \, \mathrm{d}x_1 \cdots \mathrm{d}x_n=\dfrac{n}{3}$.

    对于 $\dfrac{n(n-1)}{2}$ 个 $x_ix_j$ 项($i\neq j$),有

    \begin{align*}
        \int_0^1\,\mathrm{d}x_i\cdots\int_0^1\,\mathrm{d}x_i\cdots\int_0^1x_ix_j\,\mathrm{d}x_j\cdots\int_0^1\,\mathrm{d}x_n
         & =\int_0^1\,\mathrm{d}x_i\int_0^1x_ix_j\,\mathrm{d}x_j \\
         & =\dfrac{1}{2}\int_0^1 x_i\,\mathrm{d}x_i               \\
         & =\dfrac{1}{4}
    \end{align*}

    因此 $\int \cdots \int_{[0,1]^n} 2\sum_{i\neq j}x_ix_j \, \mathrm{d}x_1 \cdots \mathrm{d}x_n=\dfrac{n(n-1)}{4}$.

    从而总和为 $\dfrac{n}{3}+\dfrac{n(n-1)}{4}=\dfrac{n(3n+1)}{12}$.
\end{sol}

(3) $\displaystyle\int_{0}^{1}\,\mathrm{d}x_1 \int_{0}^{x_1}\,\mathrm{d}x_2 \cdots \int_{0}^{x_{n-1}} x_1 \cdots x_n \, \mathrm{d}x_n$.
\begin{sol}
    \begin{align*}
        \int_{0}^{1}\,\mathrm{d}x_1 \int_{0}^{x_1}\,\mathrm{d}x_2 \cdots \int_{0}^{x_{n-1}} x_1 \cdots x_n \, \mathrm{d}x_n
         & =\int_{0}^{1}x_1\,\mathrm{d}x_1 \int_{0}^{x_1}x_2\,\mathrm{d}x_2 \cdots \int_{0}^{x_{n-1}} x_n \, \mathrm{d}x_n                              \\
         & =\int_{0}^{1}x_1\,\mathrm{d}x_1 \int_{0}^{x_1}x_2\,\mathrm{d}x_2 \cdots \int_{0}^{x_{n-2}} \dfrac{1}{2}x_{n-1}^3 \, \mathrm{d}x_{n-1}         \\
         & =\int_{0}^{1}x_1\,\mathrm{d}x_1 \int_{0}^{x_1}x_2\,\mathrm{d}x_2 \cdots \int_{0}^{x_{n-3}} \dfrac{1}{2\times 4}x_{n-2}^5 \, \mathrm{d}x_{n-2} \\
         & =\dots                                                                                                                                       \\
         & =\int_0^1\dfrac{1}{2^{n-1}(n-1)!}x_1^{2n-1}\,\mathrm{d}x_1                                                                                    \\
         & =\dfrac{1}{2^nn!}
    \end{align*}
\end{sol}

3. 设 $f(x)$ 连续,证明:
$$\int_{0}^{a}\,\mathrm{d}x_1 \int_{0}^{x_1}\,\mathrm{d}x_2 \cdots \int_{0}^{x_{n-1}} f(x_n) \, \mathrm{d}x_n = \frac{1}{(n-1)!} \int_{0}^{a} f(t)(a - t)^{n-1} \, \mathrm{d}t.$$
\begin{pf}
    积分区域为 $0\leq x_n\leq x_{n-1}\leq\dots\leq x_1\leq a$.由于最内层被积函数仅与 $x_n$ 有关,固定 $x_n=t$,则 $t\leq x_{n-1}\leq\dots\leq x_1\leq a$.变换积分顺序,可得

    \begin{align*}
        \int_{0}^{a}\,\mathrm{d}x_1 \int_{0}^{x_1}\,\mathrm{d}x_2 \cdots \int_{0}^{x_{n-1}} f(x_n) \, \mathrm{d}x_n
         & =\int_0^a f(t)\,\mathrm{d}t\int_{t\leq x_{n-1}\leq\dots\leq x_1\leq a}\,\mathrm{d}x_1\cdots\mathrm{d}x_{n-1} \\
         & =\int_0^a f(t)\dfrac{(a-t)^{n-1}}{(n-1)!}\,\mathrm{d}t                                                        \\
         & =\dfrac{1}{(n-1)!} \int_{0}^{a} f(t)(a - t)^{n-1} \, \mathrm{d}t
    \end{align*}

    Q.E.D.
\end{pf}

4. 设 $f(x)$ 连续,证明:
$$\int_{0}^{a}\,\mathrm{d}x_1 \int_{0}^{x_1}\,\mathrm{d}x_2 \cdots \int_{0}^{x_{n-1}} f(x_1)f(x_2) \cdots f(x_n) \, \mathrm{d}x_n = \frac{1}{n!} \left( \int_{0}^{a} f(t) \, \mathrm{d}t \right)^n.$$
\begin{pf}
    \begin{align*}
         & \quad\int_{0}^{a}\,\mathrm{d}x_1 \int_{0}^{x_1}\,\mathrm{d}x_2 \cdots \int_{0}^{x_{n-1}} f(x_1)f(x_2) \cdots f(x_n) \, \mathrm{d}x_n \\
         & =\int_{0\leq x_n\leq\dots\leq x_1\leq a}f(x_1)f(x_2) \cdots f(x_n)\,\mathrm{d}x_1\cdots\mathrm{d}x_n                                 \\
         & =\dfrac{1}{n!} \left( \int_{0}^{a} f(t) \, \mathrm{d}t \right)^n
    \end{align*}

    Q.E.D.
\end{pf}

\section*{第10章综合习题}

8. 设 $a,b$ 是不全为 $0$ 的常数. 求证:
$$
    \iint\limits_{x^2+y^2 \leq 1} f(ax + by + c) \, \mathrm{d}x\mathrm{d}y = 2 \int_{-1}^{1} \sqrt{1-t^2} \, f\left(t\sqrt{a^2+b^2} + c\right) \, \mathrm{d}t.
$$
\begin{pf}
    记 $r=\sqrt{a^2+b^2}$,采用正交变换:
    $$\begin{cases}
        u=\dfrac{ax+by}{r} \\
        v=\dfrac{-bx+ay}{r}
    \end{cases},$$
    则 $u^2+v^2\leq 1$,且雅可比行列式 $\dfrac{\partial(x,y)}{\partial(u,v)}=1$.

    此时被积函数化为 $f(ur+c)$.

    因此

    \begin{align*}
        \iint_{x^2+y^2 \leq 1} f(ax + by + c) \, \mathrm{d}x\mathrm{d}y
         & =\int_{u^2+v^2\leq 1}f(ur+c)\,\mathrm{d}u\mathrm{d}v                             \\
         & =\int_{-1}^1\,\mathrm{d}u\int_{-\sqrt{1-u^2}}^{\sqrt{1-u^2}}f(ur+c)\,\mathrm{d}v \\
         & =2\int_{-1}^1\sqrt{1-u^2}f(ur+c)\,\mathrm{d}u
    \end{align*}

    令 $t=u$,则有

    \begin{align*}
        \iint_{x^2+y^2 \leq 1} f(ax + by + c) \, \mathrm{d}x\mathrm{d}y
         & =2\int_{-1}^1\sqrt{1-u^2}f(ur+c)\,\mathrm{d}u                                     \\
         & =2 \int_{-1}^{1} \sqrt{1-t^2} \, f\left(t\sqrt{a^2+b^2} + c\right) \, \mathrm{d}t
    \end{align*}

    Q.E.D.
\end{pf}

\end{document}